% from ega2-1-fr.tex
\oldpage[II]{55}

où $M$ est un $S$-module gradué, et
$\mathcal{I}=\widetilde{\mathfrak{I}}$, où $\mathfrak{I}$ est un idéal de
$A$. Pour tout $f$ homogène de $S_+$, la restriction à
$D_+(f)=\operatorname{Spec}(S_{(f)})$ du premier membre de
(\hyperref[II.3.2.4.1-fr]{3.2.4.1}) est associée à
$(\mathfrak{I}M)_{(f)}=\mathfrak{I}\cdot M_{(f)}$, et il en est de même de la
restriction du second membre, vu
(\hyperref[I.1.3.13-fr]{I, 1.3.13} et
\hyperref[I.1.6.9-fr]{I, 1.6.9}).

\begin{proposition}[3.2.5]
\label{II.3.2.5-fr}
Soit $f\in\Gamma(Y,\mathcal{S}_d)$ ($d>0$). Sur l'ensemble ouvert $X_f$, le
$(\mathcal{O}_X|X_f)$-Module
$\widetilde{\mathcal{S}(nd)}|X_f$ est canoniquement isomorphe à
$\mathcal{O}_X|X_f$ pour tout $n\in\mathbf{Z}$. En particulier, si la
$\mathcal{O}_Y$-Algèbre $\mathcal{S}$ est engendrée par $\mathcal{S}_1$
(\hyperref[II.3.1.9-fr]{3.1.9}), les $\mathcal{O}_X$-Modules
$\widetilde{\mathcal{S}(n)}$ sont inversibles pour tout $n\in\mathbf{Z}$.
\end{proposition}

\begin{proof}
En effet, pour tout ouvert affine $U$ de $Y$, on a défini dans
(\hyperref[II.2.5.7-fr]{2.5.7}) un isomorphisme canonique de
$\widetilde{\mathcal{S}(nd)}|(X_f\cap\varphi^{-1}(U))$ sur
$\mathcal{O}_X|(X_f\cap\varphi^{-1}(U))$, compte tenu de
(\hyperref[II.3.1.4-fr]{3.1.4}) (où $\varphi$ est le morphisme structural
$X\to Y$)~; il est immédiat de vérifier que ces isomorphismes sont compatibles
avec la restriction de $U$ à un ouvert affine $U'\subset U$, d'où la première
assertion. Pour établir la seconde, il suffit de remarquer que si
$\mathcal{S}$ est engendrée par $\mathcal{S}_1$, il y a un recouvrement
$(U_\alpha)$ de $Y$ par des ouverts affines tels que
$\Gamma(U_\alpha,\mathcal{S})$ soit engendrée par
$\Gamma(U_\alpha,\mathcal{S})_1=\Gamma(U_\alpha,\mathcal{S}_1)$~; on est alors
ramené à appliquer le résultat de (\hyperref[II.2.5.9-fr]{2.5.9}), la
propriété d'être inversible étant locale.
\end{proof}

On posera encore, pour tout $n\in\mathbf{Z}$,
\[
\label{II.3.2.5.1-fr}
  \mathcal{O}_X(n)=\widetilde{\mathcal{S}(n)}
\tag{3.2.5.1}
\]
et pour tout $\mathcal{O}_X$-Module $\mathcal{F}$
\[
\label{II.3.2.5.2-fr}
  \mathcal{F}(n)=\mathcal{F}\otimes_{\mathcal{O}_X}\mathcal{O}_X(n).
\tag{3.2.5.2}
\]
Il résulte aussitôt de ces définitions que pour tout ouvert $U$ de $Y$, on a
\[
  \widetilde{(\mathcal{S}|U)(n)}
  =\mathcal{O}_X(n)|f^{-1}(U)
\]
$f$ étant le morphisme structural $X\to Y$.

\begin{proposition}[3.2.6]
\label{II.3.2.6-fr}
Soient $\mathcal{M}$ et $\mathcal{N}$ deux $\mathcal{S}$-Modules gradués
quasi-cohérents, il existe un homomorphisme canonique fonctoriel (en
$\mathcal{M}$ et $\mathcal{N}$)
\[
\label{II.3.2.6.1-fr}
  \lambda:\widetilde{\mathcal{M}}
  \otimes_{\mathcal{O}_X}\widetilde{\mathcal{N}}
  \longrightarrow\widetilde{\mathcal{M}\otimes_{\mathcal{S}}\mathcal{N}}
\tag{3.2.6.1}
\]
et un homomorphisme canonique fonctoriel (en $\mathcal{M}$ et
$\mathcal{N}$)
\[
\label{II.3.2.6.2-fr}
  \mu:\widetilde{\mathcal{H}om_{\mathcal{S}}(\mathcal{M},\mathcal{N})}
  \longrightarrow
  \mathcal{H}om_{\mathcal{O}_X}
  (\widetilde{\mathcal{M}},\widetilde{\mathcal{N}}).
\tag{3.2.6.2}
\]
En outre, si $\mathcal{S}$ est engendrée par $\mathcal{S}_1$
(\hyperref[II.3.1.9-fr]{3.1.9}), $\lambda$ est un isomorphisme~; si de plus
$\mathcal{M}$ admet une présentation finie
(\hyperref[II.3.1.1-fr]{3.1.1}), $\mu$ est un isomorphisme.
\end{proposition}

\begin{proof}
Les isomorphismes $\lambda$ et $\mu$ ont été définis dans
(\hyperref[II.2.5.11.2-fr]{2.5.11.2}) et
(\hyperref[II.2.5.12.2-fr]{2.5.12.2}) lorsque $Y$ est affine~; ces définitions
étant locales se transportent aussitôt au cas général considéré ici, compte
tenu de (\hyperref[II.2.8.14-fr]{2.8.14}).
\end{proof}

\begin{corollary}[3.2.7]
\label{II.3.2.7-fr}
Si $\mathcal{S}$ est engendrée par $\mathcal{S}_1$, on a, quels que soient
$m$, $n$ dans $\mathbf{Z}$,
\[
\label{II.3.2.7.1-fr}
  \mathcal{O}_X(m)\otimes_{\mathcal{O}_X}\mathcal{O}_X(n)
  =\mathcal{O}_X(m+n)
\tag{3.2.7.1}
\]
\[
\label{II.3.2.7.2-fr}
  \mathcal{O}_X(n)=(\mathcal{O}_X(1))^{\otimes n}
\tag{3.2.7.2}
\]
à des isomorphismes canoniques près.
\end{corollary}

